\begin{ttuabstract}
Quantum processors with over a hundred qubits are now on commercial clouds, yet none completes a computation of practical value without much classical computation. This dissertation argues that the classical high-performance computing (HPC) around a quantum processing unit (QPU) is a permanent, load-bearing part of the architecture. It advances five layers of that co-designed stack. At the simulation layer, Q-GEAR transforms Qiskit circuits gate-by-gate into CUDA-Q kernels through a compact three-tensor encoding. It deploys them as containers on Slurm-managed GPU systems and simulates up to 42 qubits on 1,024 NVIDIA A100 GPUs of the Perlmutter supercomputer. Speed-ups are roughly two orders of magnitude over a 128-core CPU node and about one over existing GPU simulators. At the data layer, four learning primitives on vectorized address--data encodings are validated on superconducting and trapped-ion hardware. QPIE is a non-sequential hybrid network with parallel classical--quantum information exchange. VQT is a quantum transformer that computes attention by a vectorized quantum dot product with a learnable nonlinear encoder. ShardQ, a hardware-aware circuit cutting-and-knitting compiler, reduces encoding error by about 15\%. NNQA compiles trained polynomial networks into exact quantum arithmetic with over 99.5\% accuracy for polynomials up to degree 35.

At the algorithm layer, direct entanglement ansatz learning (DEAL) replaces black-box QAOA optimization. It uses a direct mapping from QUBO coefficients to circuit angles, a calibration-driven dynamic physical mapping and adaptive digital zero-noise extrapolation. DEAL raises ground-state success on IBM Heron processors by up to 14\%. The quantum approximate walk algorithm (QAWA) introduces a linear-depth classical-data-traceable oracle with outputs verifiable in polynomial time. At the synthesis layer, RubriQ frames fault-tolerant circuit synthesis as rubric-guided reinforcement learning of a code-generating language model, with GPU simulation inside the reward loop. RubriQ achieves 3.31$\times$ mean T-gate compression with fewer than 1\% hardware-constraint violations. A unified cost-and-error model for direct quantum simulation of nonlinear waves quantifies where measurement and reload costs overtake the classical alternative.

At the trust layer, a bounded cloud fault model and the notion of accepted logical disturbance make error-correcting encoders auditable against a co-located fault-injection adversary. A polynomial-cost seeded Clifford ensemble reduces the mean accepted logical disturbance of eight-qubit encoders from 0.150 to 0.020, an 86.7\% reduction. The fixed five-qubit code corrects every tested weight-one fault. A post-quantum cryptographic module designed to the FIPS 140-3 requirements supplies validated randomness and key-establishment services to the audit and the rest of the stack. Together these contributions show that gains at each classical layer compound, and that the path to useful quantum computation runs through the whole stack.
\end{ttuabstract}
